\section{训练系统效率：异步架构与采样瓶颈}

讲者中段明确指出，当前大规模 RL 训练的主要瓶颈已从 trainer 逐步转向 sampler。原因在于长轨迹 rollout、环境交互、grader 调用带来的吞吐压力远高于传统 batch training。

\begin{knowledgebox}{为什么异步（asynchronous）几乎是必选项}
同步模式下 sampler 与 trainer 相互等待，GPU/CPU 利用率都会下降。异步模式允许采样与训练解耦，并通过策略版本管理缓解 stale policy 带来的偏差。
\end{knowledgebox}

课程里提到 ``sampler 需要更多 GPU'' 的现象，对很多初学者是反直觉的。传统认知认为训练最耗算力，但在 agent RL 中，环境+rollout 可能成为主耗时路径，尤其当每个轨迹都需要多轮工具调用时。

\begin{table}[H]
\centering
\begin{tabular}{p{2.8cm}p{3.8cm}p{3.9cm}}
\toprule
\textbf{组件} & \textbf{主要耗时} & \textbf{优化抓手} \\
\midrule
Sampler / Roller & 长轨迹生成、环境调用、grader 请求 & 并行 rollout、缓存、早停策略 \\
Trainer & 反向传播、参数同步 & mixed precision、梯度累计、LoRA \\
Grader Service & 可执行验证与模型判分 & 分层调度、异步队列、批处理 \\
Env Sandbox & 测试运行与状态隔离 & 轻量容器、快照回滚、冷热池 \\
\bottomrule
\end{tabular}
\caption{Agent RL 中各模块效率瓶颈与优化方向}
\end{table}

\begin{importantbox}{效率与质量不是零和，但需要结构化平衡}
讲者强调团队常使用 ``更快但次优'' 的 baseline 做大规模对比实验，先加速创新迭代，再在最终版本上用最优配置拉满质量。这种阶段化策略能把研发吞吐提升到 2x--3x。
\end{importantbox}

\subsection{本章小结}
异步架构是工业 Agent RL 的基础设施。真正的优化目标不是单次最高分，而是单位时间内可复现、可迭代的稳定增益。

